% Options for packages loaded elsewhere
% Options for packages loaded elsewhere
\PassOptionsToPackage{unicode}{hyperref}
\PassOptionsToPackage{hyphens}{url}
\PassOptionsToPackage{dvipsnames,svgnames,x11names}{xcolor}
%
\documentclass[
  letterpaper,
  DIV=11,
  numbers=noendperiod]{scrartcl}
\usepackage{xcolor}
\usepackage{amsmath,amssymb}
\setcounter{secnumdepth}{-\maxdimen} % remove section numbering
\usepackage{iftex}
\ifPDFTeX
  \usepackage[T1]{fontenc}
  \usepackage[utf8]{inputenc}
  \usepackage{textcomp} % provide euro and other symbols
\else % if luatex or xetex
  \usepackage{unicode-math} % this also loads fontspec
  \defaultfontfeatures{Scale=MatchLowercase}
  \defaultfontfeatures[\rmfamily]{Ligatures=TeX,Scale=1}
\fi
\usepackage{lmodern}
\ifPDFTeX\else
  % xetex/luatex font selection
\fi
% Use upquote if available, for straight quotes in verbatim environments
\IfFileExists{upquote.sty}{\usepackage{upquote}}{}
\IfFileExists{microtype.sty}{% use microtype if available
  \usepackage[]{microtype}
  \UseMicrotypeSet[protrusion]{basicmath} % disable protrusion for tt fonts
}{}
\makeatletter
\@ifundefined{KOMAClassName}{% if non-KOMA class
  \IfFileExists{parskip.sty}{%
    \usepackage{parskip}
  }{% else
    \setlength{\parindent}{0pt}
    \setlength{\parskip}{6pt plus 2pt minus 1pt}}
}{% if KOMA class
  \KOMAoptions{parskip=half}}
\makeatother
% Make \paragraph and \subparagraph free-standing
\makeatletter
\ifx\paragraph\undefined\else
  \let\oldparagraph\paragraph
  \renewcommand{\paragraph}{
    \@ifstar
      \xxxParagraphStar
      \xxxParagraphNoStar
  }
  \newcommand{\xxxParagraphStar}[1]{\oldparagraph*{#1}\mbox{}}
  \newcommand{\xxxParagraphNoStar}[1]{\oldparagraph{#1}\mbox{}}
\fi
\ifx\subparagraph\undefined\else
  \let\oldsubparagraph\subparagraph
  \renewcommand{\subparagraph}{
    \@ifstar
      \xxxSubParagraphStar
      \xxxSubParagraphNoStar
  }
  \newcommand{\xxxSubParagraphStar}[1]{\oldsubparagraph*{#1}\mbox{}}
  \newcommand{\xxxSubParagraphNoStar}[1]{\oldsubparagraph{#1}\mbox{}}
\fi
\makeatother


\usepackage{longtable,booktabs,array}
\usepackage{calc} % for calculating minipage widths
% Correct order of tables after \paragraph or \subparagraph
\usepackage{etoolbox}
\makeatletter
\patchcmd\longtable{\par}{\if@noskipsec\mbox{}\fi\par}{}{}
\makeatother
% Allow footnotes in longtable head/foot
\IfFileExists{footnotehyper.sty}{\usepackage{footnotehyper}}{\usepackage{footnote}}
\makesavenoteenv{longtable}
\usepackage{graphicx}
\makeatletter
\newsavebox\pandoc@box
\newcommand*\pandocbounded[1]{% scales image to fit in text height/width
  \sbox\pandoc@box{#1}%
  \Gscale@div\@tempa{\textheight}{\dimexpr\ht\pandoc@box+\dp\pandoc@box\relax}%
  \Gscale@div\@tempb{\linewidth}{\wd\pandoc@box}%
  \ifdim\@tempb\p@<\@tempa\p@\let\@tempa\@tempb\fi% select the smaller of both
  \ifdim\@tempa\p@<\p@\scalebox{\@tempa}{\usebox\pandoc@box}%
  \else\usebox{\pandoc@box}%
  \fi%
}
% Set default figure placement to htbp
\def\fps@figure{htbp}
\makeatother





\setlength{\emergencystretch}{3em} % prevent overfull lines

\providecommand{\tightlist}{%
  \setlength{\itemsep}{0pt}\setlength{\parskip}{0pt}}



 


\KOMAoption{captions}{tableheading}
\makeatletter
\@ifpackageloaded{caption}{}{\usepackage{caption}}
\AtBeginDocument{%
\ifdefined\contentsname
  \renewcommand*\contentsname{Table of contents}
\else
  \newcommand\contentsname{Table of contents}
\fi
\ifdefined\listfigurename
  \renewcommand*\listfigurename{List of Figures}
\else
  \newcommand\listfigurename{List of Figures}
\fi
\ifdefined\listtablename
  \renewcommand*\listtablename{List of Tables}
\else
  \newcommand\listtablename{List of Tables}
\fi
\ifdefined\figurename
  \renewcommand*\figurename{Figure}
\else
  \newcommand\figurename{Figure}
\fi
\ifdefined\tablename
  \renewcommand*\tablename{Table}
\else
  \newcommand\tablename{Table}
\fi
}
\@ifpackageloaded{float}{}{\usepackage{float}}
\floatstyle{ruled}
\@ifundefined{c@chapter}{\newfloat{codelisting}{h}{lop}}{\newfloat{codelisting}{h}{lop}[chapter]}
\floatname{codelisting}{Listing}
\newcommand*\listoflistings{\listof{codelisting}{List of Listings}}
\makeatother
\makeatletter
\makeatother
\makeatletter
\@ifpackageloaded{caption}{}{\usepackage{caption}}
\@ifpackageloaded{subcaption}{}{\usepackage{subcaption}}
\makeatother
\usepackage{bookmark}
\IfFileExists{xurl.sty}{\usepackage{xurl}}{} % add URL line breaks if available
\urlstyle{same}
\hypersetup{
  pdftitle={Problem Set 3},
  colorlinks=true,
  linkcolor={blue},
  filecolor={Maroon},
  citecolor={Blue},
  urlcolor={Blue},
  pdfcreator={LaTeX via pandoc}}


\title{Problem Set 3}
\author{}
\date{}
\begin{document}
\maketitle


\begin{enumerate}
\def\labelenumi{\arabic{enumi}.}
\item
  \textbf{The Secret to Longevity}. Scientists are interested in whether
  the energy costs involved in reproduction affect longevity\footnote{Question
    from Problem 3.2 in Oehlert textbook, original data from Hanley and
    Shapiro (1994), ``Sexual activity and the lifespan of male
    fruitflies: a dataset that gets attention,'' \emph{Journal of
    Statistics Education 2}.}. In this experiment, 125 male fruit flies
  were divided at random into five sets of 25. In group 1, the males
  were kept by themselves. In groups 3 and 5, the males were supplied
  with one or eight receptive virgin female fruit flies per day,
  respectively. In groups 2 and 4, the males were supplied with one or
  eight unreceptive (pregnant) female fruit flies per day, respectively.
  Other than the number and type of companions, the males were treated
  identically. The longevity of the flies was observed.

  You can access the data as a csv file at:
  \url{https://stat158.berkeley.edu/spring-2026/data/fruit-flies/fruit-flies.csv}

  \begin{enumerate}
  \def\labelenumii{\alph{enumii}.}
  \item
    Visualize the data and interpret what you see. Be sure to comment
    both on the center of the distributions and their spread.
  \item
    \emph{ANOVA ``by hand''}. Without using any of the built-in commands
    (\texttt{aov()}, \texttt{anova()}) create a one-way ANOVA table
    minus the p-values. You're welcome to use either a base R or
    \texttt{dplyr} approach. It is a bit tedious, however it's useful to
    cobble together each of the statistics in an ANOVA table just once
    in your lives. You can check your answers with the output of
    \texttt{aov()} shown below.
  \end{enumerate}

\begin{verbatim}
             Df Sum Sq Mean Sq F value Pr(>F)
factor(trt)   4  11939  2984.8   13.61       
Residuals   120  26314   219.3               
\end{verbatim}

  \begin{enumerate}
  \def\labelenumii{\alph{enumii}.}
  \setcounter{enumii}{2}
  \item
    State your conclusion about the scientific question using
    Model-based inference. Be clear about the data generation model your
    approach relies upon and the null hypothesis that you're evaluating.
    To calculate the p-value, use the \texttt{pf()} function and then
    verify it using the \texttt{aov()} function\footnote{Recall all
      named probability distributions have a CDF in R with the prefix
      \texttt{p}. The F-distribution is no exception, so you can use
      \texttt{pf()} to calculate the p-value for your ANOVA table. The
      arguments to \texttt{pf()} are the F-statistic, the degrees of
      freedom for the numerator, and the degrees of freedom for the
      denominator. Recall a test using the F statistic is a right-tailed
      test. The standard anova output is available with
      \texttt{summary(aov(y\ \textasciitilde{}\ x,\ data))} (be sure
      your \texttt{x} is a factor vector).}.
  \item
    State your conclusion about the scientific question using
    Randomization-based Inference. Be clear about the data generation
    model your approach relies upon and the null hypothesis that you're
    evaluating. For your test statistic, use the F-statistic that you
    built the code for from part a above.
  \item
    How do the results of the model-based and randomization-based
    inference compare? Do they lead to the same conclusion? Do they rely
    on the same assumptions? Which do you think is more appropriate for
    this problem?
  \end{enumerate}
\end{enumerate}

\begin{enumerate}
\def\labelenumi{\arabic{enumi}.}
\setcounter{enumi}{1}
\item
  \textbf{More Comparisons for Longevity}. In the experiment regarding
  the fruit flies, the ANOVA analysis only tests whether there were any
  differences between the group means. Suppose there was further
  interest in whether the 1 pregnant group was different from the 1
  virgin group (for questions a and b, restrict the data set to these
  two groups).

  \begin{enumerate}
  \def\labelenumii{\alph{enumii}.}
  \item
    Perform a randomization test to compare these two groups using the
    t-statistic. Assume the two groups have the same variance (revisit
    your notes on the t-test and the pooled estimate of the SE).
  \item
    Repeat the randomization test but this time use an F-statistic. How
    do the results of this analysis this compare to the results in part
    (a)?
  \item
    Calculate an ANOVA table for this comparison and report the value of
    \(\sqrt{MS_{residuals}}\), which is an estimate of \(\sigma\). How
    does that compare to the estimate of the standard deviation that you
    used for the t-statistic (\(s_{p}\))?
  \end{enumerate}

  You are welcome to use built-in R commands for this question, if
  appropriate.
\end{enumerate}

\newpage{}

More questions to come next week!




\end{document}
